% Figure: electrical-to-luminous efficiency waterfall for the white LED
% luminaire case study. Each bar is the running fraction of input electrical
% power surviving one conversion step, ending at the wall-plug luminous
% efficacy expressed as a fraction of the warm-white radiometric ceiling.
% Plain TikZ rectangles; same idiom as the solar-cell waterfall.
\begin{figure}[htbp]
\centering
\begin{tikzpicture}[font=\scriptsize]
\def\sy{0.06}
\draw[enggray, thin] (0,0) -- (12.7,0);
\foreach \v in {0,20,40,60,80,100} {
  \draw[enggray, very thin] (-0.1,\v*\sy) -- (12.7,\v*\sy);
  \node[enggray, anchor=east, font=\tiny] at (-0.15,\v*\sy) {\v};
}
\node[enggray, anchor=south, rotate=90, font=\scriptsize]
  at (-0.9,50*\sy) {fraction of input electrical power (\%)};
% 1 electrical input 100
\filldraw[engblue]   (0.3,0)  rectangle (1.6,100*\sy);
\node[anchor=south, font=\tiny] at (0.95,100*\sy) {100};
\node[anchor=north, font=\tiny, align=center] at (0.95,0) {electrical\\input};
% 2 after electrical efficiency (forward-voltage / series-resistance), ~85
\filldraw[engamber]  (2.0,0)  rectangle (3.3,85*\sy);
\node[anchor=south, font=\tiny] at (2.65,85*\sy) {85};
\node[anchor=north, font=\tiny, align=center] at (2.65,0) {after\\$V_f$ loss};
% 3 after internal quantum efficiency ~70 (0.82 IQE)
\filldraw[engorange] (3.7,0)  rectangle (5.0,70*\sy);
\node[anchor=south, font=\tiny] at (4.35,70*\sy) {70};
\node[anchor=north, font=\tiny, align=center] at (4.35,0) {after\\IQE};
% 4 after extraction efficiency ~49 (0.70 extraction)
\filldraw[engred]    (5.4,0)  rectangle (6.7,49*\sy);
\node[anchor=south, font=\tiny] at (6.05,49*\sy) {49};
\node[anchor=north, font=\tiny, align=center] at (6.05,0) {after\\extraction};
% 5 after phosphor down-conversion (Stokes + quantum) ~39
\filldraw[englime]   (7.1,0)  rectangle (8.4,39*\sy);
\node[anchor=south, font=\tiny] at (7.75,39*\sy) {39};
\node[anchor=north, font=\tiny, align=center] at (7.75,0) {after\\phosphor};
% 6 luminous efficiency of the radiation (spectrum vs eye), running ~28
\filldraw[enggray]   (8.8,0)  rectangle (10.1,28*\sy);
\node[anchor=south, font=\tiny] at (9.45,28*\sy) {28};
\node[anchor=north, font=\tiny, align=center] at (9.45,0) {luminous\\of radiation};
% 7 after thermal derating at junction temperature ~24
\filldraw[engblue]   (10.5,0) rectangle (11.8,24*\sy);
\node[anchor=south, font=\tiny] at (11.15,24*\sy) {24};
\node[anchor=north, font=\tiny, align=center] at (11.15,0) {after\\derating};
\end{tikzpicture}
\caption[White-LED luminaire efficiency waterfall]{Where the power goes in a
phosphor-converted white-LED luminaire, as a running fraction of the input
electrical power. The forward-voltage and series-resistance drop, the internal
radiative quantum efficiency, the photon-extraction efficiency from the
high-index die, the phosphor down-conversion (Stokes shift plus the
phosphor's own quantum yield), the luminous efficiency of the resulting
warm-white spectrum against the eye's response, and the thermal derating at the
operating junction temperature each subtract a separate factor. The last bar at
about $24\,\%$ of input power corresponds to a wall-plug luminous efficacy near
$160\,\si{\lumen\per\watt}$ against the $683\,\si{\lumen\per\watt}$ monochromatic
ceiling. Each step is a distinct physical loss with a distinct remedy, and the
case study carries them in order.}
\label{fig:vol04ch12:led-efficiency-waterfall}
\end{figure}
